\begin{table}[t]
\tbl{\label{tab:pc:activity}}
{\begin{tcolorbox}[tabularx={p{3cm}|p{2cm}|p{6.3cm}|},enhanced,fonttitle=\bfseries\large,fontupper=\normalsize\sffamily,
colback=yellow!10!white,colframe=red!50!black,colbacktitle=blue!30!white,
coltitle=black,title=Activities for Posets of Orbits]
\hline
Command & Arguments & Purpose \\
\hline
\hline
\verb'-report'  &  options   &  Produce a latex report. For options, see Table~\ref{tab:pc:report:options}. Requires \verb'-orbiter_path' option from Section~\ref{sec:session}. \\
\hline
\verb'-export_level_to_cpp'  &  i   &  Export level i to cpp file. \\
\hline
\verb'-export_history_to_cpp'  &  i   &  Export history up to level i to cpp file. \\
\hline
\verb'-write_tree'  &  options   &  Create the depth first tree from the poset of orbits. A tree file is created. Uses the given draw options.\\
\hline
\verb'-find_node_by_stabilizer_order'  &  k   &  Search for all nodes whose stabilizer has order k. \\
\hline
\verb'-draw_poset'  &  options   &  Draw the poset of orbits. Uses the given draw options.\\
\hline
\verb'-draw_full_poset'  &   options  &  Draw the poset of orbits in full. Uses the given draw options. \\
\hline
\verb'-plesken_ring'  &     &  Compute the Plesken matrix of the poset of orbits. See~\cite{Plesken82}. This command returns an object of type Plesken\_ring. \\
\hline
\verb'-print_data_structure'  &     &  Print the data structure of the poset of orbits. \\
\hline
\verb'-list'  &     &  List the nodes in the poset of orbits.  \\
\hline
\verb'-list_all'  &     &  List all nodes in the poset of orbits.  \\
\hline
\verb'-table_of_nodes'  &     &  Produce a table of all nodes in the poset of orbits. \\
\hline
\verb'-make_relations_with_flag_orbits'  &     &  Produce a bitmap drawing of the neighboring relations in the poset with flag orbits. \\
\hline
\verb'-level_summary_csv'  &     &  Write a summary of number of orbits at each level to a csv file.  \\
\hline
\verb'-orbit_reps_csv'  &     &  Write orbit representatives to a csv file. \\
\hline
\verb'-node_label_is_group_order'  &     &  When drawing the poset of orbits, display the group order in the orbit nodes. \\
\hline
\verb'-node_label_is_element'  &     &  When drawing the poset of orbits, display the element rank in the orbit nodes. \\
\hline
\verb'-show_orbit_decomposition'  &     &  Show the orbits of the stabilizers on the orbit representatives (in the induced action on the set). \\
\hline
\verb'-show_stab'  &     &  Show the stabilizers. \\
\hline
\verb'-save_stab'  &     &  Save the stabilizer generators. \\
\hline
\verb'-show_whole_orbits'  &     &  Show the whole orbits. \\
\hline
\verb'-export_schreier_trees'  &     &  Export all Schreier trees. \\
\hline
\verb'-draw_schreier_trees'  &     &  Draw all Schreier trees.  \\
\hline
\verb'-test_multi_edge_in_' \verb'decomposition_matrix'  &     &  \\
\hline
\verb'-recognize'  &  set   &  Recognizes the given set in the poset of orbits. This option can be repeated.\\
\hline
\verb'-pair_relations_within_orbit'  &  orbit   &  Computes the pairwise join of elements in the given orbit and identifies the orbit in the poset. Two matrices are written. In the first matrix, the $ij$-entry is the index of the orbit containing the join of orbit element $i$ and orbit element $j$. In the second matrix, the $ij$-entry is the orbit index of the join of orbit element $i$ and orbit element $j$ in its respective orbit. \\
\hline
\verb'-export_orbits_long'  &     &  Exports all orbits in long form. For each orbit, each orbit element is listed as a set.\\
\hline
\end{tcolorbox}}
\end{table}
